\section*{Part IX: The Founder's Adjacency Decision System}
\addcontentsline{toc}{section}{Part IX: The Founder's Adjacency Decision System}
\label{sec:part9_founders_adjacency_decision_system}

\subsection{The 1-Page Adjacency Thesis}
\label{sec:one_page_adjacency_thesis}

A credible expansion proposal should state: current core advantage; new buyer and job; existing substitutes; how customers actually pay; what transfers from the core; what new capability must be built; expected incremental contribution; working-capital and founder-hour needs; the largest tail risk; legal permissions; asset that may accumulate; how rivals respond; the smallest revealing experiment; and a predefined stopping rule. Add the best alternative use of the same resources.

Write the thesis before commissioning a large forecast. Its purpose is not to persuade a committee but to expose assumptions that have been hidden in the attractiveness of the story. If the responsible owner cannot describe the new buyer's purchase process or the new service's fulfillment obligations, the proposal is not ready for funding.

\subsection{A 2-Stage Decision: Hard Gates Then Comparative Ranking}
\label{sec:two_stage_decision_gates_then_ranking}

1st apply non-compensatory gates: lawful and ethical operation; credible buyer problem; feasible route to payment; acceptable downside to existing obligations; availability of essential expertise; and adequate cash. Failing any gate means reject, delay, or redesign. Then compare surviving candidates on economic contribution, capability transfer, option value, reversibility, complexity, and capital intensity. Scores are discussion aids, not permission for a weak project to pass a fatal test.

Make uncertainty visible. An apparently attractive project's score might depend heavily on an untested customer-conversion assumption. Mark it as provisional and design an experiment around that assumption rather than disguising the problem through a precise weighted total.

\subsection{The Adjacency Distance Worksheet}
\label{sec:adjacency_distance_worksheet}

\begin{table*}[t]
\centering
\small
\caption{The Adjacency Distance Worksheet: Comparing Existing Capabilities Against Expansion Requirements}
\label{tab:adjacency_distance_worksheet}
\begin{tabularx}{\textwidth}{p{2.6cm} p{2.6cm} p{2.6cm} p{2.8cm} p{2.8cm} X}
\toprule
\textbf{Dimension} & \textbf{Existing Model} & \textbf{Proposed Model} & \textbf{Transferable Asset} & \textbf{New Learning or Cost} & \textbf{Evidence Status} \\
\midrule
Buyer \& Job & Working software engineers & University engineering deans & Rigorous evaluation rubrics & Academic curriculum alignment & Unvalidated hypothesis \\
\addlinespace[2pt]
Purchasing Process & Self-funded credit card / UPI & Institutional procurement & Transparent syllabus specs & 90-day purchase order cycle & Initial dean interview \\
\addlinespace[2pt]
Distribution Channel & Inbound technical essays & Direct enterprise outreach & Demonstrable case studies & Cold enterprise outreach skill & 2 inbound requests \\
\addlinespace[2pt]
Delivery \& Quality & Weekend interactive studio & Semester-long weekly lab & Automated grading tests & Live classroom TA coordination & Pilot lab tested \\
\addlinespace[2pt]
Talent \& Management & Founder-led facilitation & Delegated adjunct reviewers & Rubric calibration manuals & Recruited TA grading oversight & 1 TA trained \\
\addlinespace[2pt]
Regulation \& Law & Standard private contract & University MOU / vendor MSA & Standard terms of service & Intellectual property assignment & Draft terms reviewed \\
\addlinespace[2pt]
Cash Flow Timing & Prepaid before cohort start & Net-60 post-delivery milestone & None (requires float) & Delayed invoice receivables & Bank reserve required \\
\bottomrule
\end{tabularx}
\end{table*}

After filling the table, identify the \textbf{2 hardest transitions}. Ask what would be required to prove each without taking the whole risk. The worksheet is strongest when accompanied by named accounts, sample contracts, actual quotes from service providers, and observed customer decisions.

\subsection{The Contribution and Capacity Worksheet}
\label{sec:contribution_and_capacity_worksheet}

For each candidate record realistic paid units, realized revenue net of discounts and refunds, incremental acquisition, direct delivery labor, reviewer or support hours, technology, travel, legal obligations, receivables timing, 1-off build cost, staffing steps, and residual operating commitments. Compute contribution and the lowest projected cash balance under adverse demand and payment schedules. Compare with the core alternative using bottleneck hours and working capital.

Do not mix cash and accounting timing. An advance for a semester program may increase the bank balance while creating a delivery liability. A large institutional invoice may show revenue before cash arrives. If the new activity cannot be modeled at this level, it has not earned an irreversible commitment.

\subsection{Real-Option Experiment Design}
\label{sec:real_option_experiment_design}

Define the important uncertainty and what information would actually change the investment choice. Choose the smallest faithful experiment. Set exposure limits in money, time, customer commitment, and reputation. Predefine evidence thresholds and what happens after a positive, negative, or inconclusive result. Record the value of waiting, including competitive timing risk and the availability of partners.

A sensible assessment-pilot option might test paid employer willingness to use a competency rubric before developing a large platform. A cross-border pilot could test real paid participation and support costs before establishing a foreign company. A research option could test feasibility before acquiring infrastructure. The experiment is useful only if it changes a future decision.

\subsection{A Board-Level Integration or Acquisition Memo}
\label{sec:board_level_acquisition_memo}

A proposed acquisition should include strategic rationale, stand-alone normalized financials, credible synergies with owners and probabilities, full purchase and integration costs, working-capital needs, key-person exposure, customer contracts, IP and privacy rights, culture, downside scenarios, financing terms, decision rights, and walk-away price. Require a clear comparison with licensing, hiring, alliances, or internal build. Ask who would manage the acquired business the following Monday.

A founder should not use complex financing to hide weak economics. When the organization has no integration bench, the price should reflect that risk, or the deal should not be done. Protect existing customer commitments before adding a new balance-sheet obligation.

\subsection{The Portfolio Review Cadence}
\label{sec:portfolio_review_cadence}

Monthly, review new proposals and commitments that threaten current capacity. Quarterly, assess the portfolio's cash contribution, founder-hour allocation, quality, customer concentration, shared-asset creation, and contingent liabilities. Annually, reconsider whether the business boundaries still match its mission and changing customer demand. New research may justify moving closer to a capability; new technology may justify exiting an obsolete service.

Assign owners to decisions, not just dashboards. A recurring review that never terminates weak projects can institutionalize drift. Document why each new activity continues to deserve scarce resources compared with a strong alternative.

\subsection{The Expansion Pre-Mortem}
\label{sec:expansion_pre_mortem}

Before proceeding, imagine that after 18 months the adjacency has damaged the core. Possible explanations: underestimated procurement delay; weak customer fit; founder attention diverted; unqualified delivery staff; incompatible regulatory obligations; overstated synergy; customer concentration; system complexity; brand confusion; supplier dependence; or cash shortfall. Ask which warning signal would appear \textit{early enough} to act.

Then run the reverse pre-mortem: imagine an apparently successful 1st pilot generates dangerous overconfidence. What might be special about its customer, buyer relationship, timing, or pricing? A single enthusiastic flagship institution is not repeatability. State how many comparable deployments and what delegated delivery would change the judgment.

\subsection{A Disciplined Answer to ``Should We Launch This?''}
\label{sec:disciplined_answer_should_we_launch}

The best response is not always a binary yes or no. It may be: ``We will sell a narrow pilot to 1 authorized buyer, preserve the core delivery schedule, spend no more than this amount, and reconsider after independently measured results.'' Or: ``This is a separate consulting business, not an adjacency, and we are not resourced to operate it.'' Or: ``Our existing offer has larger unexploited upside; we should improve it 1st.'' Or: ``Technological change is undermining the core, so the option cost of waiting is high and controlled experimentation is urgent.''

Good strategy is not a permanent refusal to diversify. It is the ability to \textbf{identify the terms under which diversification becomes rational} and to change those terms when evidence changes.
